\chapter{The Architecture of CoALA}
\label{ch:coala-architecture}

\begin{learningbox}
By the end of this chapter, you should be able to read CoALA as one connected architecture rather than as a list of terms. You will be able to identify its three organizing dimensions, explain the relationship between the structural view and the decision-cycle view, distinguish the LLM from the larger agent, and trace how information, actions, and control fit together. The chapter deliberately previews memory and action categories without replacing the detailed treatment in Chapters 3--5.
\end{learningbox}

\section{From a collection of capabilities to an architecture}

Imagine an agent that can call a language model, retain information from earlier tasks, retrieve documents, invoke tools, and follow a program. We can list those capabilities, perhaps even identify the software components that implement them. But a list does not yet explain how they form one agent. Which information is active when a decision is made? How does the model use that information? Which operations can change internal information, and which can affect the environment? What decides which operation happens next?

Return to Product X, our book-created software incident scenario. The payment service is returning HTTP 500 errors. The agent may need the current incident report, recent deployment information, facts about the service, records of earlier incidents, and procedures for checking production systems. It may need to inspect logs, retrieve a runbook, reason about a possible cause, ask an operator a question, or run a health check. Those are not interchangeable activities. Some read retained information; some process information already available; some write new information; others interact with the running service or with a person.

The important question is not simply whether each capability exists. It is how each capability relates to the others and how the system coordinates their use. A component can be technically present but disconnected from the decision process. A tool can be registered but never selected. A memory store can contain useful records that are not retrieved at the right time. A model can produce a plausible instruction that no procedure actually executes. To understand the agent, we need to understand its organization.

In this chapter, we use the word \term{architecture} in a practical teaching sense: a map of the major elements of a system, the responsibilities assigned to them, and the relationships through which information and control move. This is an explanatory description for the reader, not a new formal definition claimed to appear verbatim in the CoALA paper. It is also not a deployment diagram. It does not tell us which cloud provider hosts a service, how many processes run, or which database product stores a record. Nor is it a class diagram that prescribes source-code types. It is a conceptual map that helps us ask what the agent contains and how its parts cooperate.

\begin{tipbox}
When reading an architecture diagram, do not stop at the boxes. Ask what each box is responsible for, what information enters and leaves it, and which part of the system determines when the interaction occurs.
\end{tipbox}

CoALA offers a common way to organize these questions. It describes a language agent along three dimensions: \term{information storage}, \term{action space}, and \term{decision-making procedure}. The LLM is central, but it is one component inside a larger cognitive architecture.\CoALACite{Section 4 and Figure 4}

\section{The three dimensions that organize CoALA}

The three dimensions can be introduced with three questions:

\begin{center}
\begin{tabularx}{\textwidth}{>{\bfseries}p{0.25\textwidth} X}
\toprule
Architectural dimension & Question it answers \\
\midrule
Information storage & What information does the agent maintain, and where is it available? \\
Action space & What can the agent do with its internal information, and how can it interact with the outside world? \\
Decision-making procedure & How does the agent choose and coordinate actions over time? \\
\bottomrule
\end{tabularx}
\end{center}

These are three views of one system, not three independent products that happen to be installed side by side. Information storage describes what the agent can draw on. The action space describes how the agent can use internal information and interact with its environment. The decision-making procedure coordinates the use of those capabilities. A description that omits one of the three dimensions leaves an important architectural question unanswered.

\subsection{Information storage: what information is maintained?}

A language model call receives an input and produces an output. In the simplest interaction, information supplied for one call is not automatically maintained as a structured state for the next call. A \gls{languageagent} that operates across multiple steps may need to preserve information between calls and decision cycles. CoALA therefore distinguishes \gls{workingmemory} from \gls{longtermmemory}. Long-term memory, in turn, includes episodic, semantic, and procedural forms.\CoALACite{Section 4.1}

For now, we need only the architectural contrast. Working memory is the active information for the current decision cycle. Long-term memory retains information that can be used beyond that immediate cycle. The long-term categories distinguish experience, knowledge, and how-to information. Chapter 3 will examine their meanings, boundaries, and examples in detail. Here, their purpose is to show that an architecture needs more than one undifferentiated store called “memory.”

\subsection{Action space: what can the agent do?}

An agent must do more than contain information. It must have operations for using it. CoALA distinguishes external actions from internal actions. External actions interact with an environment through grounding. Internal actions interact with the agent's own memory. The paper further describes three kinds of internal action: retrieval, reasoning, and learning.\CoALACite{Sections 4.2--4.5 and Figure 5}

At the level of this chapter, the distinctions are simple. Retrieval reads from long-term memory. Reasoning uses the LLM to process information in working memory. Learning writes information to long-term memory. Grounding executes external actions and processes environmental feedback. Chapters 4 and 5 will examine the operational details. For the architectural map, notice that “action” is broader than “call a tool”: some actions operate on internal information, while others affect or obtain information from the external environment.

\subsection{Decision-making procedure: how are actions coordinated?}

If an agent has information and operations, we still need to know how it determines what to do. CoALA describes decision-making as a repeated process involving planning and execution. During a decision cycle, the agent can use retrieval and reasoning to support planning. Planning proposes and evaluates candidate learning or grounding actions; a selected action is then executed. An observation may follow, and the next cycle begins.\CoALACite{Section 4.6 and Figure 4B}

This high-level description is enough to understand why the decision-making procedure is a separate architectural dimension. The action space tells us which operations are available. The decision procedure tells us how the agent coordinates them. We will separate planning, evaluation, selection, and execution in Chapters 6--8; this chapter establishes only the overall relationship.

\begin{takeawaybox}[The architecture in one sentence]
CoALA organizes a language agent by asking what information it maintains, what internal and external actions it can perform, and what decision procedure coordinates those actions in a repeated interaction with the world.
\end{takeawaybox}

\section{The structural map: components and relationships}

The three questions can be read as a set of responsibilities, but CoALA also shows how the elements relate. Its structural view places the LLM and agent code alongside working memory, long-term memory, the decision procedure, and interfaces to the external environment. The connections matter: the agent code constructs prompts and parses model outputs; memory supplies and receives information; the decision procedure invokes the action procedures; grounding carries actions and observations across the boundary between the agent and its environment.\CoALACite{Figure 4A and Section 4}

\begin{figure}[htbp]
\centering
\resizebox{\textwidth}{!}{%
\begin{tikzpicture}[
  box/.style={draw=BookBlue, rounded corners=2pt, thick, fill=NoteBg, minimum width=2.2cm, minimum height=0.85cm, align=center, font=\sffamily\small},
  memory/.style={draw=BookBlue, rounded corners=2pt, thick, fill=LearnBg, minimum width=2.35cm, minimum height=0.75cm, align=center, font=\sffamily\small},
  action/.style={draw=BookMid, rounded corners=2pt, fill=TipBg, minimum width=1.8cm, minimum height=0.68cm, align=center, font=\sffamily\small},
  flow/.style={-{Stealth[length=2.5mm]}, thick, draw=BookBlue},
  back/.style={-{Stealth[length=2.5mm]}, thick, draw=BookMid, dashed}
]
\draw[rounded corners=5pt, thick, draw=BookBlue] (-0.35,-3.45) rectangle (10.65,2.5);
\node[font=\sffamily\bfseries\color{BookBlue},anchor=west] at (-0.1,2.22) {Cognitive language agent};
\node[box] (llm) at (1.3,1.2) {LLM};
\node[box] (code) at (4.5,1.2) {Agent code};
\node[memory] (wm) at (7.8,1.2) {Working memory};
\node[box] (decision) at (1.7,-1.0) {Decision procedure};
\node[action] (retrieval) at (4.7,-0.35) {Retrieval};
\node[action] (reasoning) at (6.8,-0.35) {Reasoning};
\node[action] (learning) at (8.9,-0.35) {Learning};
\node[action] (grounding) at (9.0,-1.4) {Grounding};
\node[memory] (ltm) at (6.8,-2.35) {Long-term memory};
\node[box] (env) at (12.3,-2.35) {External environment};
\draw[flow] (code.west) -- (llm.east);
\node[font=\scriptsize,align=center] at (2.9,1.95) {prompt construction / output parsing};
\draw[flow] (wm.west) -- (code.east);
\node[font=\scriptsize] at (6.15,1.95) {active context};
\draw[flow] (decision) -- (code);
\draw[flow] (code) -- (retrieval);
\draw[flow] (code) -- (reasoning);
\draw[flow] (code) -- (learning);
\draw[flow] (retrieval) -- (ltm);
\draw[flow] (learning) -- (ltm);
\draw[back] (ltm) -- (wm);
\draw[flow] (code.south east) .. controls (6.0,-0.2) and (7.6,-1.4) .. (grounding.west);
\draw[flow] (grounding.east) -- (env.north west);
\draw[back] (env.west) .. controls (10.7,-3.05) and (9.1,-3.15) .. node[below,font=\scriptsize]{observations} (wm.south);
\end{tikzpicture}%
}
\caption{An original teaching redraw of the structural relationships in CoALA's Figure 4A. The drawing is schematic: arrows identify information or control relationships, not a required software topology.}
\label{fig:coala-structural-map}
\end{figure}

The figure is a teaching redraw, not a reproduction of the paper's artwork. Its purpose is to help us read the relationships. It does not prescribe that an implementation must use separate machines, processes, classes, or databases for every box. Nor should every arrow be mistaken for a network call. The diagram expresses roles and interactions at the conceptual level.

\subsection{The LLM and agent code}

The LLM is a core computational component of a \gls{cognitivelanguageagent}. It supplies language-based processing that can support reasoning and generate outputs used by the agent. But CoALA does not identify the LLM with the whole agent. The architecture also includes agent code: explicit procedures that construct model inputs, parse outputs, implement actions, and execute the decision-making process.\CoALACite{Section 4 and Figure 4A}

This distinction is easier to see if we separate two questions. First, what capabilities are available through the model? The LLM's learned parameters encode implicit knowledge that can influence the output it generates. Second, what procedures does the surrounding system explicitly execute? Agent code defines how prompts are built, how model outputs are interpreted, how memory is accessed, how grounding actions are performed, and how the decision process proceeds.

The paper also describes \gls{proceduralmemory} as having implicit and explicit forms: knowledge stored in the LLM weights and knowledge written in the agent's code. These are not identical storage mechanisms, even though both can contribute to how the agent behaves. Chapter 3 will examine this distinction as part of the wider memory architecture. For now, remember that an LLM call does not itself supply every operation needed to run the agent.

\subsection{Working memory as the active information hub}

Working memory occupies a central position because it maintains active information for the current decision cycle. CoALA describes it as a data structure that persists across LLM calls. The LLM input is synthesized from a subset of working memory, such as a prompt template populated with relevant variables. The model's output is parsed and stored back into variables in working memory, where it can be used to execute a corresponding action. Working memory also interacts with long-term memory and grounding interfaces.\CoALACite{Section 4.1}

That description gives us a useful correction to a common assumption: working memory is not simply another name for the current prompt. A prompt is an input representation prepared for one model invocation. Working memory is the larger active information structure from which that input can be assembled and into which results can be written. The current prompt may contain only the information relevant to a specific call.

Consider Product X. The active working information might include the incident goal, the latest observation, the current hypothesis, and the result of the most recent check. A particular model call may need only the latest error pattern and a request to compare two hypotheses. Another call may need a different subset, perhaps the result of a deployment check. The architecture can maintain more active information than is placed in any one prompt.

\begin{notebox}
The term “working memory” is used here in CoALA's architectural sense. Do not equate it automatically with a model vendor's context window, conversation transcript, cache, or a specific in-memory data structure. An implementation may connect these ideas, but the conceptual role is broader than any one mechanism.
\end{notebox}

\subsection{Long-term memory at the architectural level}

CoALA differentiates long-term memory into \gls{episodicmemory}, \gls{semanticmemory}, and procedural forms. The categories answer different questions: what happened in prior experiences, what the agent knows about the world and itself, and how the agent performs actions and decision procedures. The paper discusses both implicit procedural knowledge in model weights and explicit procedural knowledge in code.\CoALACite{Section 4.1}

This chapter does not attempt to define every category in full. It is enough to understand why the architecture distinguishes them. If the agent keeps an incident history, stores facts about the payment service, and contains executable procedures for checking the service, those are not the same kind of information. The design question is not simply “where is the memory database?” It is “what kind of information is this, what role does it play, and how does the agent access or update it?”

Chapter 3 will explore the memory categories, their relationships, and the implications of treating memory as organized information rather than as one transcript. Here the only point is architectural: long-term memory is connected to working memory through internal actions, and the decision procedure determines when those actions are used.

\subsection{The environment and the two-way connection}

The agent is not isolated inside its memory modules. \gls{grounding} procedures execute external actions and process environmental feedback into working memory. The CoALA paper describes physical environments, dialogue with humans or other agents, and digital environments as examples of external settings.\CoALACite{Section 4.2}

In Product X, a digital environment may expose service health, log records, deployment history, and database metrics. A dialogue environment may include an operator who can clarify whether a change was authorized. A physical environment is less central to this incident, but the same conceptual boundary applies to a robot receiving sensor data and issuing motor commands. The environment differs, while the architectural role of grounding remains recognizable.

The connection is two-way. Actions go from the agent toward the environment; observations come back. An action that is merely proposed in text has not yet crossed that boundary. A grounding procedure must execute the interaction and process the resulting feedback. That feedback can then become active information in working memory and influence a later decision.

\subsection{Reading the connections, not just the boxes}

An architecture diagram can be misleading when a reader treats it as an inventory. The value of the structural view lies in the relationships.

\begin{itemize}
  \item The decision procedure invokes explicit procedures in agent code.
  \item Agent code interacts with the LLM by preparing input and parsing output.
  \item Working memory provides active information for the current cycle and receives results.
  \item Retrieval and learning connect working memory with long-term memory through reads and writes.
  \item Grounding connects the agent with the external environment through actions and observations.
\end{itemize}

A single component may participate in several relationships. The LLM can help reason over working information, but it does not by itself decide how every operation is executed. Working memory may be the active information hub, but it is not the only memory store. Agent code implements procedures, but it is the decision procedure that coordinates how they are used over time. The environment contributes observations and receives actions, but it does not belong to the agent's internal memory architecture.

\section{Memory's place in the architecture}

It is useful to pause on memory because the structural diagram places it at the centre of several connections. A reader who thinks of memory as a single database may miss why CoALA separates working memory from long-term memory and then differentiates long-term memory into subtypes. The purpose of this section is not to teach those subtypes in full; it is to clarify their place in the architecture.

Working memory supports the current decision cycle. It holds relevant observations, active goals, knowledge retrieved from long-term memory, and results generated during reasoning. Long-term memory allows the agent to maintain information beyond the current cycle. Retrieval can bring relevant information into working memory; learning can write information from working memory into long-term memory. These connections make memory part of the agent's ongoing activity rather than a passive archive.

The categories also show that “remembering” is not one operation. Reading an earlier incident, retaining a fact about a service, and using an executable procedure are different roles. An architecture that labels all of them merely “memory” may hide important design differences. Conversely, the names of the categories do not require separate physical databases. The conceptual distinctions can be represented in different ways in software, provided the implementation makes their responsibilities clear.

\begin{warningbox}
Do not infer from the diagram that CoALA mandates a particular storage product, schema, vector database, file format, persistence layer, or retention policy. Those are implementation decisions. CoALA provides the conceptual memory distinctions; engineering choices must be justified separately.
\end{warningbox}

Chapter 3 will unpack working memory and the long-term memory categories. The current chapter only establishes the links: memory makes information available, actions read or write it, and the decision procedure determines when those operations are used.

\section{The action space: internal capabilities and external interaction}

The structural view also connects the agent to two broad kinds of action. External actions interact with an environment through grounding. Internal actions interact with internal memories. The internal category is further divided into retrieval, reasoning, and learning.\CoALACite{Sections 4.2--4.5 and Figure 5}

\subsection{Internal actions}

The three internal action categories are distinguished by what information they access and how they change it.

\begin{description}
  \item[Retrieval] reads information from long-term memory. Its role is to make relevant retained information available for the current decision process.
  \item[Reasoning] uses the LLM to process information in working memory and generate new information there. It reads from and writes to working memory.
  \item[Learning] writes information to long-term memory. In the CoALA account, learning is not limited to changing neural-network weights; it can include storing information in memory modules.
\end{description}

These definitions are deliberately brief. Chapter 4 will explain retrieval, reasoning, and learning as operational categories, including their directions of information movement and examples. This chapter only needs the distinction to make sense of the map.

\subsection{External action: grounding}

\gls{grounding} procedures execute external actions and process feedback from the environment into working memory. The paper groups external interaction into physical environments, dialogue with humans or other agents, and digital environments. A robot may receive visual and tactile observations and issue motor commands. A dialogue agent may ask a person a question or communicate with another agent. A digital agent may navigate a website, interact with an API, play a game, or execute code.\CoALACite{Section 4.2}

In each case, the action is more than a sentence describing an intention. It is an operation performed through an interface to the external world. The system must interpret the result as an observation that can inform later decisions. The agent's environment therefore forms part of the architecture's boundary: it is outside the agent's internal memory, but it is connected to the agent through grounding.

\subsection{Why planning is not the same as the entire action space}

Figure 5 in the CoALA paper provides a compact view of these categories. Reasoning and retrieval are internal actions that support planning. Learning is also an internal action, but it is not placed within the planning bracket in the figure. Grounding is external. This is an important distinction: planning uses certain internal actions to determine which learning or grounding action to select; it is not a synonym for every action available to the agent.\CoALACite{Figure 5}

\begin{figure}[htbp]
\centering
\begin{tikzpicture}[
  act/.style={draw=BookBlue, thick, fill=TipBg, minimum width=1.25cm, minimum height=0.8cm, align=center, font=\sffamily\small},
  group/.style={draw=BookBlue, thick, rounded corners=2pt}
]
\node[act] (reason) at (0,0) {Reasoning};
\node[act] (retrieval) at (1.55,0) {Retrieval};
\node[act] (learning) at (3.1,0) {Learning};
\node[act] (grounding) at (4.65,0) {Grounding};
% Grouping brackets are deliberately outside the action nodes so that the
% categories remain readable at print size.
\draw[thick,draw=BookBlue] (-0.68,0.72) -- (3.78,0.72);
\draw[thick,draw=BookBlue] (-0.68,0.72) -- (-0.68,0.58);
\draw[thick,draw=BookBlue] (3.78,0.72) -- (3.78,0.58);
\draw[thick,draw=BookBlue] (3.98,0.72) -- (5.32,0.72);
\draw[thick,draw=BookBlue] (3.98,0.72) -- (3.98,0.58);
\draw[thick,draw=BookBlue] (5.32,0.72) -- (5.32,0.58);
\node[font=\sffamily\bfseries\small,BookBlue] at (1.55,1.0) {Internal};
\node[font=\sffamily\bfseries\small,BookBlue] at (4.65,1.0) {External};
\draw[thick,draw=BookBlue] (-0.68,-0.72) -- (2.32,-0.72);
\draw[thick,draw=BookBlue] (-0.68,-0.72) -- (-0.68,-0.58);
\draw[thick,draw=BookBlue] (2.32,-0.72) -- (2.32,-0.58);
\node[font=\sffamily\small,BookBlue] at (0.82,-1.04) {Planning support};
\end{tikzpicture}
\caption{An original redraw of CoALA's action-space distinction. Reasoning and retrieval support planning; learning remains an internal action; grounding is external.}
\label{fig:coala-action-space}
\end{figure}

The distinction helps avoid a common overgeneralization. An agent's external tool calls are not its whole action space, because internal memory operations also matter. Conversely, not every internal operation is itself a planning step. Chapters 4 and 5 will examine the actions in detail; Chapters 6--8 will explain how planning works with them.

\section{The decision procedure: how the architecture operates over time}

A structural map tells us which elements exist and how they relate, but it does not fully show their order over time. For that, CoALA presents a \gls{decisionprocedure} view. The paper compares the decision procedure to a program's main procedure: a top-level routine that runs repeatedly, accepts new perceptual input, and calls action procedures in response.\CoALACite{Section 4.6}

\subsection{The top-level or main procedure}

Think of an ordinary software application. A program may contain functions for reading input, validating data, writing a file, and returning a result. The existence of those functions does not explain when they run. A top-level procedure determines the control flow: which functions are called, in what context, and what happens after they return.

CoALA's decision procedure plays a comparable organizing role for the agent. It coordinates action procedures and repeats the interaction with the environment. This comparison is a teaching analogy based on the paper's description; it does not mean that every implementation must contain a method literally named `main`, nor that the full decision process is a single fixed code listing.

\subsection{Planning and execution at a glance}

The paper's temporal view places planning before execution. During planning, retrieval and reasoning can help the agent propose and evaluate candidate learning or grounding actions. Selection chooses an action. Execution then carries out the selected action. An observation may be made, after which the decision cycle begins again.\CoALACite{Section 4.6 and Figure 4B}

\begin{figure}[htbp]
\centering
\begin{tikzpicture}[
  node distance=0.55cm,
  stage/.style={draw=BookBlue, thick, rounded corners=2pt, fill=TipBg, minimum width=2.2cm, minimum height=0.62cm, align=center, font=\sffamily\small},
  container/.style={draw=BookBlue, thick, rounded corners=4pt, inner sep=8pt},
  flow/.style={-{Stealth[length=2.4mm]}, thick, draw=BookBlue}
]
\node[stage] (obs) {Observation};
\node[stage,below=of obs] (proposal) {Proposal};
\node[stage,below=of proposal] (evaluation) {Evaluation};
\node[stage,below=of evaluation] (selection) {Selection};
\node[stage,below=0.8cm of selection] (execution) {Execution};
\draw[container] (-1.6,-3.3) rectangle (1.6,-0.4);
\node[font=\sffamily\bfseries\color{BookBlue}] at (0,-0.16) {Planning};
\draw[flow] (obs) -- (proposal);
\draw[flow] (proposal) -- (evaluation);
\draw[flow] (evaluation) -- (selection);
\draw[flow] (selection) -- (execution);
\draw[flow] (execution.east) -- ++(1.0,0) |- (obs.east) node[pos=0.28,right,font=\scriptsize,align=left]{next cycle / feedback};
\end{tikzpicture}
\caption{An original teaching redraw of CoALA's decision-cycle view. The stages are shown in the explanatory order used by the figure; the detailed chapters will examine interleaving and iteration.}
\label{fig:coala-decision-cycle}
\end{figure}

It is important not to read this high-level picture as a complete low-level algorithm. The figure helps us see the broad relation between planning and execution. It does not mean that the agent must perform one isolated reasoning call, then one isolated retrieval, then one proposal, then one evaluation, with no possibility of returning to an earlier subtask. The details of how planning's sub-stages relate, and how they may be revisited, are the subject of Chapters 6--8.

\subsection{A cycle is a useful view, not a rigid script}

The word \term{cycle} emphasizes repetition. It reminds us that an agent may receive new observations after an action and must decide what to do in light of them. The cycle also shows that the decision procedure does not merely select external tool calls. It can select a learning action that updates long-term memory as well as a grounding action that interacts with the environment.

At the same time, the diagram should not be mistaken for a universal, mandatory sequence that every CoALA agent must follow identically. CoALA is a conceptual framework. The paper describes planning in terms of proposal, evaluation, and selection, but the architecture is intended to organize varied agents and decision procedures. The detailed account in Chapters 6--8 will preserve the specific roles of these stages without imposing a rigid one-pass pipeline where the paper does not.

\section{A complete architectural reading of Product X}

We can now return to the payment incident and read it through the architecture. The following scenario is created for this book. Its purpose is to apply CoALA's distinctions, not to claim that the paper presents this incident or prescribes the specific operational steps.

The payment service is returning HTTP 500 errors. The agent's goal is to investigate the incident and provide a supported recommendation. It begins with an observation from the monitoring system: the error rate increased shortly after deployment 4.7. The agent does not yet know whether the deployment caused the failure.

\subsection{Information storage}

Working memory contains the active incident goal, the latest observations, relevant facts already in hand, and the current state of the investigation. Long-term memory may contain records of earlier incidents, facts about the payment service, and procedures for checking its health. The architecture distinguishes the role of active information from information retained for future use.

\subsection{Action space}

The agent may retrieve previous incidents or service facts from long-term memory. It may reason over the deployment timestamp and the error increase to formulate candidate explanations. It may learn by writing a new experience or fact to long-term memory when the procedure calls for it. It may ground by requesting a health check, reading logs through an external interface, or asking an operator for clarification.

These operations are different even when they are triggered by the same investigation. Retrieving a runbook does not itself inspect the service. Reasoning about a possible database failure does not query the database. A grounding procedure performs the external interaction and processes the result as an observation.

\subsection{Decision-making procedure}

The decision procedure coordinates what happens next. It can use retrieval and reasoning to support planning. It can propose candidate actions such as inspect database connectivity or compare the recent deployment with the error timeline. It can evaluate the candidates and select one for execution. Once the selected operation is carried out, an observation may update working memory. The next cycle can then take the new information into account.

The exact investigation strategy, authorization process, rollback policy, and operational safety controls are not defined by the CoALA architecture map itself. They would be engineering choices for a real incident-response system. CoALA helps us describe the cognitive organization of the agent; it does not automatically make the system safe to execute production changes.

\begin{takeawaybox}[What the example demonstrates]
A single incident can involve every architectural dimension. The agent maintains active and retained information, uses internal and external actions, and relies on a decision procedure to coordinate them. The example makes the architecture concrete without adding new CoALA mechanisms.
\end{takeawaybox}

\section{One framework, different agent configurations}

A framework is useful only if it can describe more than one narrowly constructed system. CoALA is intended to organize a range of language agents, not to require one fixed collection of components or one identical decision strategy. The paper applies the framework to different systems and uses it to compare their memory, actions, and decision-making characteristics.\CoALACite{Section 4 and Table 2}

This means that two agents can both be understood through CoALA while differing substantially in what they implement. One may have a simple memory arrangement and a small action space. Another may maintain several forms of long-term memory and use a more elaborate decision procedure. One agent may primarily interact through dialogue; another may act in a digital environment. The common framework provides the vocabulary for describing those differences.

The purpose of a common vocabulary is not to erase variation. It is to make variation easier to see. Instead of saying only that one system is “more agentic” or “uses more tools,” we can ask: What information does it maintain? What internal operations does it support? What external environment does it interact with? How does its decision procedure select and execute operations? Where are its capabilities limited?

This is also why a framework should not be confused with a product specification. CoALA gives us a way to describe and reason about language-agent organization. A particular application still needs requirements, domain constraints, interfaces, reliability properties, safety policies, and evaluation methods. Those are not supplied automatically by drawing the architecture.

\section{Common misreadings and architectural boundaries}

Several misunderstandings can arise when the architecture is reduced to a few memorable labels. It is worth correcting them before moving to the detailed chapters.

\begin{enumerate}
  \item \textbf{The LLM is the agent.} The LLM is central, but the architecture also includes agent code, memory, action procedures, and a decision procedure.
  \item \textbf{Working memory is just the prompt.} The prompt is synthesized from a relevant subset of working memory. Working memory is the active information structure that can persist across model calls.
  \item \textbf{All actions are tool calls.} CoALA includes internal memory actions as well as external grounding actions. Retrieval, reasoning, and learning are not interchangeable with external tool execution.
  \item \textbf{Planning is the whole action space.} Planning uses reasoning and retrieval to support candidate generation and evaluation. Learning is also an internal action, and grounding is an external action.
  \item \textbf{The diagram is a deployment plan.} The boxes are conceptual roles. They do not require one service, process, class, or database per box.
  \item \textbf{The decision cycle is a rigid one-pass script.} The cycle is a high-level view of repeated decision-making. Chapters 6--8 will examine the roles and revisiting of planning sub-stages carefully.
  \item \textbf{A CoALA description proves production readiness.} The framework is not a substitute for authorization, testing, security, reliability, observability, or domain validation.\index{architecture!cognitive} \index{memory!working} \index{memory!long-term} \index{grounding} \index{decision procedure}
\end{enumerate}

\begin{warningbox}
Do not import mechanisms from Fan and Lan into the explanation of the original CoALA framework. Part II will study that later mechanism-level analysis separately. Likewise, do not attribute implementation choices made in Part III to CoALA unless the original paper explicitly supports the claim.
\end{warningbox}

\section{How the remaining CoALA chapters build on this map}

This chapter has provided the map. The next chapters will zoom in on individual regions without losing the relationships between them.

\begin{itemize}
  \item \textbf{Chapter 3 --- Understanding Memory.} Working memory, episodic memory, semantic memory, procedural memory, and the boundaries between them.
  \item \textbf{Chapter 4 --- Internal Actions: Retrieval, Reasoning and Learning.} The internal operations, their inputs and outputs, and how they affect memory.
  \item \textbf{Chapter 5 --- External Actions and Grounding.} External environments, grounding procedures, observations, and action examples.
  \item \textbf{Chapter 6 --- Planning: From a Problem to Candidate Actions.} The proposal stage and the generation of candidate actions.
  \item \textbf{Chapter 7 --- Evaluation and Selection.} Candidate evaluation, values, selection methods, and the possibility of rejecting candidates.
  \item \textbf{Chapter 8 --- Interleaving, Iteration and the Complete CoALA Decision Process.} How the planning stages can be revisited and interleaved, and how the complete decision process operates.
\end{itemize}

The order matters pedagogically. We first need the whole map, then we can inspect its components and processes in more detail. The chapters will return to the same Product X example so that the reader can see how a single incident draws on memory, actions, and decision-making. The example remains a teaching aid, not a source claim.

\section{Chapter recap}

\begin{chapterrecap}
CoALA organizes language agents through three dimensions: information storage, action space, and decision-making procedure. The LLM is a core component within a larger architecture. Agent code implements procedures that connect the LLM to working memory, long-term memory, and the external environment. Internal actions include retrieval, reasoning, and learning; external actions use grounding. The decision procedure coordinates planning and execution in a repeated cycle. The structural diagram explains the components and their relationships, while the temporal diagram explains how the agent operates over time. Both views describe the same architecture from different angles.
\end{chapterrecap}

\begin{reviewquestions}{Review questions}
\begin{enumerate}
  \item Why is a list of capabilities not enough to explain an agent's architecture?
  \item What are the three dimensions CoALA uses to organize language agents?
  \item Why is working memory broader than the prompt supplied to one LLM call?
  \item How do internal actions differ from external grounding actions?
  \item Why are reasoning and retrieval associated with planning in the action-space diagram?
  \item What does the structural view show that the decision-cycle view does not, and vice versa?
  \item Why should the decision-cycle diagram not be read as a rigid one-pass script?
  \item Which decisions in the Product X example are illustrations rather than requirements prescribed by CoALA?
\end{enumerate}
\end{reviewquestions}

\section*{Source and example note}
The theoretical account in this chapter is grounded in \textsc{[CoALA]} only, principally Section 4 and Figures 4--5 of Sumers et al. The diagrams are original teaching redraws of the paper's conceptual relationships. Product X and all incident-specific details are examples created for this book. This chapter does not use Fan and Lan to reinterpret CoALA and does not prescribe a software implementation.
